
\chapter{M002 --- Primitive Surface Construction}

\section{Stage 3 --- Implementation Review (Phase 3)}

\subsection{Reviewed Transformation}

Construction of the exact primitive kernel/B\'ezout triplet
$(u,v,w)$.

\subsection{Mathematical Input}

\begin{itemize}
\item Primitive Miller covector.
\item Verified primitive reciprocal basis.
\end{itemize}

\subsection{Mathematical Output}

A certified integer triplet $(u,v,w)$ where

\[
m\cdot u = 0,\qquad
m\cdot v = 0,\qquad
m\cdot w = 1,
\]

with $(u,v)$ spanning the primitive integer kernel.

\subsection{Algorithmic Role}

This phase constructs the exact integer objects required for all
subsequent geometric operations. Unlike floating-point surface
generation, all calculations remain in the integer lattice.

\subsection{Classification}

\begin{itemize}
\item Scientific transformation.
\item Exact integer-lattice construction.
\item No representation change.
\item No gauge operation.
\end{itemize}

\subsection{Relationship to M001}

This phase reuses the mathematical ideas established in M001:

\begin{itemize}
\item Primitive kernel construction.
\item Bézout witness construction.
\item Exact unimodular reasoning.
\end{itemize}

\subsection{Open Questions for Stage 4}

\begin{enumerate}
\item Does the implementation construct a primitive kernel basis?
\item Is the Bézout witness mathematically certified?
\item Are all integer invariants preserved throughout the construction?
\end{enumerate}
